\documentclass[11pt,a4paper]{article}

% -- Core packages --
\usepackage{amsmath,amssymb}
\usepackage{geometry}
\usepackage{graphicx}
\usepackage{booktabs}
\usepackage{multirow}
\usepackage{hyperref}
\usepackage{xcolor}
\usepackage{listings}
\usepackage{enumitem}
\usepackage{tcolorbox}
\usepackage{textcomp}

% -- Layout --
\geometry{margin=2.5cm}

\hypersetup{
  colorlinks=true,
  linkcolor=blue!60!black,
  urlcolor=blue!70!black
}

% -- Python / shell listings style --
\definecolor{codebg}{HTML}{F5F5F5}
\definecolor{codeframe}{HTML}{CCCCCC}
\definecolor{codegreen}{rgb}{0.0,0.45,0.0}
\definecolor{codegray}{rgb}{0.5,0.5,0.5}
\definecolor{codepurple}{rgb}{0.45,0.0,0.5}
\definecolor{codeblue}{rgb}{0.0,0.2,0.8}

\lstdefinestyle{pystyle}{
  backgroundcolor=\color{codebg},
  frame=single,
  rulecolor=\color{codeframe},
  basicstyle=\small\ttfamily,
  keywordstyle=\color{codeblue}\bfseries,
  commentstyle=\color{codegreen}\itshape,
  stringstyle=\color{codepurple},
  numberstyle=\tiny\color{codegray},
  breaklines=true,
  breakatwhitespace=false,
  showstringspaces=false,
  tabsize=4,
  language=Python,
  morekeywords={True,False,None},
  xleftmargin=6pt,
  xrightmargin=6pt,
}

\lstset{
  style=pystyle,
  literate={µ}{{\textmu}}1
           {·}{{\textperiodcentered}}1
           {—}{{---}}1
           {→}{{$\to$}}1
           {³}{{$^3$}}1
}

% -- Info box --
\tcbuselibrary{skins,breakable}
\newtcolorbox{infobox}[1][]{
  enhanced,
  breakable,
  colback=blue!5!white,
  colframe=blue!45!black,
  fonttitle=\bfseries,
  title=#1,
  left=6pt, right=6pt, top=4pt, bottom=4pt,
}
\newtcolorbox{warningbox}[1][]{
  enhanced,
  breakable,
  colback=orange!5!white,
  colframe=orange!70!black,
  fonttitle=\bfseries,
  title=#1,
  left=6pt, right=6pt, top=4pt, bottom=4pt,
}

% -- Title --
\title{%
  \textbf{FRED 2.0}\\[8pt]
  \large User Quickstart Manual\\[4pt]
  \normalsize FRED-ROD and FRED-OX --- a step-by-step guide through the example scripts
}
\author{FRED Development Team\\
        \small Paul Scherrer Institut / Advanced Nuclear System group}
\date{\today}

% ===========================================================================
\begin{document}
% ===========================================================================

\maketitle
\tableofcontents
\newpage

% ===========================================================================
\section{Introduction}
% ===========================================================================

FRED 2.0 is a platform for developing fuel-performance codes.  Core physics
--- heat conduction and thermo-elastic stress-strain, including
pellet-cladding gap closing/reopening --- lives in a shared platform layer.
Applications are built on top of that layer:

\begin{description}
  \item[\textbf{FRED-ROD}] solves thermal conduction and stress-strain only.
        No irradiation behaviour is modelled; users supply their own material
        correlations (including, if desired, their own base-irradiation
        effects).  This decoupling is deliberate: it keeps FRED-ROD suitable
        for controlled experimental/verification workflows.
  \item[\textbf{FRED-OX}] uses the same equations as FRED-ROD and adds
        oxide-fuel-specific physics on top: empirical base-irradiation
        correlations for fuel/cladding properties, and fission-gas
        production, which feeds back into gas pressure and gap conductance.
  \item[\textbf{FRED-M-Na}] uses the same equations as FRED-ROD and adds
        metallic-fuel-specific physics (U-Pu-Zr in sodium), including many
        correlations ported from SAS4A.  See
        \texttt{examples/fred\_m\_na/} and \texttt{doc/physics\_fred\_m\_na.md};
        Section~\ref{sec:mna-pointer} of this manual gives a short pointer to
        its examples.
\end{description}

\textbf{This manual covers FRED-ROD (Part~\ref{part:rod}) and FRED-OX
(Part~\ref{part:ox}) in depth}, walking through every example script under
\texttt{examples/fred\_rod/} and \texttt{examples/fred\_ox/} in a recommended
learning order, from the simplest case to the most complex.  Each part opens
with its own ``recommended example order'' list.

\medskip
\begin{infobox}[Prerequisites]
  All examples assume the relevant FRED Python module has been compiled and
  placed in \texttt{fred\_platform/build/}.  Add that directory to your
  Python path:
  \begin{lstlisting}[language=bash, style=pystyle]
# from the fred_platform root
make pymodule            # builds build/fred_rod.*.so, fred_ox.*.so, ...
export PYTHONPATH=$PWD/build:$PYTHONPATH
  \end{lstlisting}
  Each script also inserts the build path automatically with
  \texttt{sys.path.insert(0, \ldots)}, so you can run them directly from
  their own directory without setting \texttt{PYTHONPATH}.
\end{infobox}

\subsection{Common workflow}

Every FRED-ROD or FRED-OX script follows the same four-step pattern:

\begin{enumerate}
  \item \textbf{Geometry} --- set the node counts, radii, axial slice
        heights, roughness, and gas-plenum volume; call \texttt{build()}.
  \item \textbf{Materials} --- choose or define fuel, cladding, and gap
        material objects.
  \item \textbf{Solver setup} --- create a \texttt{FredRodSolver} (ROD) or
        \texttt{FredOxSolver} (OX), set physics toggles, boundary conditions,
        tolerances.
  \item \textbf{Run and post-process} --- call \texttt{run(tend, dtout)};
        retrieve results with the accessor methods.
\end{enumerate}

\subsection{Unit conventions}

\begin{center}
\begin{tabular}{lll}
\toprule
Quantity & Unit & Notes \\
\midrule
Temperature        & K        & Always Kelvin; subtract 273.15 for \textdegree C \\
Time               & s        & Seconds \\
Length / radius    & m        & Metres \\
Power density      & W/m\textsuperscript{3} & Volumetric heat generation \\
Pressure / stress  & MPa      & Megapascals \\
\bottomrule
\end{tabular}
\end{center}

% ===========================================================================
\part{FRED-ROD}
\label{part:rod}
% ===========================================================================

FRED-ROD is the thermo-mechanical application of the FRED 2.0 platform.  It
solves the coupled heat-conduction and thermo-elastic stress-strain equations
for a cylindrical fuel rod using the SUNDIALS IDA differential-algebraic
solver.  No irradiation behaviour is modelled; users add their own material
models (including, if desired, their own base-irradiation correlations).

\section{Recommended Example Order}
\label{sec:rod-order}

The example scripts under \texttt{examples/fred\_rod/} are best worked
through in the following order, progressing from the simplest case to the
most complex:

\begin{enumerate}
  \item \textbf{\texttt{heat\_conduction\_only}} --- heat equation alone; no mechanics.
  \item \textbf{\texttt{stress\_strain\_only}} --- elastic mechanics alone; temperatures fixed.
  \item \textbf{\texttt{heat\_conduction\_stress\_strain}} --- coupled thermo-elastic run.
  \item \textbf{\texttt{gapclose\_reopen}} --- gap close and reopen under power cycling.
  \item \textbf{\texttt{hot\_start\_demo}} --- skipping the transient heat-up with a
        pseudo-time hot start, and when \emph{not} to use it.
  \item \textbf{\texttt{hdf5\_output}} --- writing and reading results in HDF5 format.
  \item \textbf{\texttt{restart\_snapshot}} --- fault-recovery checkpoints and physical-state snapshots.
  \item \textbf{\texttt{python\_materials}} --- defining custom material correlations in Python.
  \item \textbf{\texttt{advanced\_prototyping}} --- the Python$\to$C++ prototyping loop for gap materials.
\end{enumerate}

\begin{infobox}[Why this order?]
  Examples 1--3 build up the physics one piece at a time (thermal only,
  mechanical only, then coupled) using the same dummy-material rod, so you
  can isolate solver behaviour before adding complexity.  Example 4
  introduces the gap-contact event logic that most real rods exercise.
  Example 5 (hot start) is easiest to understand \emph{after} Example 3,
  since it directly replaces the manual power-ramp workaround used there.
  Examples 6--7 are about persistence and I/O, independent of the underlying
  physics, so they can be read once you are comfortable with any of the
  earlier runs. Examples 8--9 are for users who want to extend FRED-ROD with
  new material correlations, and assume familiarity with Examples 1--3.
\end{infobox}

% ===========================================================================
\section{Example 1: Heat Conduction Only}
\label{sec:hco}
% ===========================================================================

\textbf{File:} \texttt{examples/fred\_rod/heat\_conduction\_only/run.py}

\medskip
This is the simplest possible run: heat conduction is enabled and the
stress-strain solver is disabled.  All mechanical state variables are pinned
to zero; the only degrees of freedom are the nodal temperatures.

\textbf{Scenario:}
\begin{itemize}
  \item Solid dummy fuel pellet and dummy cladding (constant properties).
  \item Constant volumetric power $q''' = 10^8$~W/m$^3$ applied at $t=0$.
  \item Constant coolant temperature $T_\text{cool} = 700$~K.
  \item Simulated for $t_\text{end} = 10$~s (about 30 thermal time constants).
\end{itemize}

\paragraph{Geometry setup.}

\begin{lstlisting}
g = fred.FuelRodGeometry()
g.nf   = 3          # 3 radial fuel nodes
g.nc   = 2          # 2 radial cladding nodes
g.nz   = 1          # 1 axial layer
g.rfi0 = 0.0        # solid pellet (no central void)
g.rfo0 = 4.2e-3     # outer fuel radius = 4.2 mm
g.rci0 = 4.3e-3     # inner cladding radius (100 µm as-fabricated gap)
g.rco0 = 4.9e-3     # outer cladding radius = 4.9 mm
g.dz0  = [0.10]     # one 10 cm axial slice
g.vgp  = 1.0e-6     # 1 cm^3 gas plenum
g.ruff = 5.0e-6     # 5 µm fuel surface roughness
g.rufc = 5.0e-6     # 5 µm cladding inner roughness
g.build()
\end{lstlisting}

The \texttt{nf} and \texttt{nc} fields set the number of radial nodes in the
fuel and cladding respectively.  With \texttt{nf=3, nc=2} there are 5 unknown
temperatures per axial layer.  The as-fabricated gap width is
$r_\text{ci0} - r_\text{fo0} = 100$~µm.

\paragraph{Physics toggles and boundary conditions.}

\begin{lstlisting}
fuel = fred.DummyFuelPellet()    # k=10 W/mK, rho=10000 kg/m3, cp=100 J/kgK
clad = fred.DummyCladding()
gap  = fred.DummyGapMaterial()   # h_gap(T) = 500 + 2.5*T  [W/(m2·K)]

solver = fred.FredRodSolver(g, fuel, clad, gap)
solver.set_enable_heat_conduction(True)   # thermal ODE active
solver.set_enable_stress_strain(False)    # mechanics frozen at zero

T0 = 700.0   # K
solver.set_initial_temperature(T0)
solver.set_coolant_temperature([0.0, 1e6], [T0, T0])    # constant at T0
solver.set_power_density_history([0.0, 1e6], [1.0e8, 1.0e8])  # step on at t=0

solver.set_coolant_pressure(5.0)     # MPa (required but not used here)
solver.set_initial_gas_pressure(0.1) # MPa
solver.set_tolerances(1e-6, 1e-8)
\end{lstlisting}

\texttt{set\_coolant\_temperature} prescribes the outer cladding surface
temperature (Dirichlet boundary condition).  \texttt{set\_power\_density\_history}
provides the volumetric heat source; values are linearly interpolated between
the supplied time points.

\paragraph{Analytical check.}

For a solid cylinder the steady-state fuel centreline temperature is:
\begin{equation}
  T_\text{cl} = T_\text{cool} + q''' R_f^2 / (4k_f)
              + q''' R_f^2 \ln(r_\text{ci}/r_\text{fo}) / k_\text{gap-eff}
              + \ldots
\end{equation}
With the constant-property dummy materials the expected centreline rise above
700~K is approximately $\Delta T \approx 137$~K, giving $T_\text{cl} \approx 837$~K.
The example prints a comparison table against the legacy FRED Fortran code.

\paragraph{Post-processing.}

\begin{lstlisting}
times  = solver.time_points()          # [n_steps] s
T_arr  = solver.temperatures()         # shape (n_steps, nz*(nf+nc))
T_peak = solver.peak_fuel_temperature()  # [n_steps] K

# Temperature at the innermost fuel node at the final time step:
T_centre_final = T_arr[-1, 0]

# Temperature at the outer cladding surface:
T_clad_outer_final = T_arr[-1, nf + nc - 1]
\end{lstlisting}

The \texttt{temperatures()} array has shape \texttt{(n\_steps, nz*(nf+nc))}.
For \texttt{nz=1} it collapses to \texttt{(n\_steps, nf+nc)}.  Nodes are
ordered fuel inner$\to$outer, then cladding inner$\to$outer.

Calling \texttt{set\_enable\_stress\_strain(False)} reduces the DAE to a
pure ODE system (all thermal equations are differential).  The solver runs
faster and the IDA initialisation is simpler.  Use this mode whenever
mechanics are not needed.


% ===========================================================================
\section{Example 2: Stress-Strain Only}
\label{sec:sso}
% ===========================================================================

\textbf{File:} \texttt{examples/fred\_rod/stress\_strain\_only/run.py}

\medskip
Heat conduction is disabled; all temperatures are frozen at the initial
value \texttt{T0 = 293.15 K} (the reference temperature for thermal expansion,
so all thermal strains are zero).  Only the thermo-elastic mechanics are
active.

\textbf{Scenario:}
\begin{itemize}
  \item Coolant pressure ramp from 5~MPa to 0.1~MPa over 100~s.
  \item Internal gas pressure constant at 0.1~MPa.
  \item As the external compression decreases, the cladding expands outward.
\end{itemize}

\paragraph{Physics and boundary conditions.}

\begin{lstlisting}
solver.set_enable_heat_conduction(False)   # temperatures frozen at T0
solver.set_enable_stress_strain(True)

T0 = 293.15   # K  — reference temperature (zero thermal expansion)
solver.set_initial_temperature(T0)

# Coolant pressure ramps from 5 MPa down to 0.1 MPa over 100 s
solver.set_coolant_pressure_history([0.0, 100.0], [5.0, 0.1])
solver.set_initial_gas_pressure(0.1)   # MPa
\end{lstlisting}

When heat conduction is disabled, the thermal ODE is replaced by an
algebraic constraint that pins every node temperature to
\texttt{T0}.  The stress-strain system is then quasi-static: at each
output step FRED-ROD solves the mechanical equilibrium equations using
the current (frozen) temperature distribution and the instantaneous
pressure boundary conditions.

\paragraph{Expected steady-state stress.}

The thin-wall approximation for the cladding hoop stress under external
pressure $p_o$ and internal pressure $p_i$ is:
\begin{equation}
  \sigma_\theta \approx \frac{p_i r_i - p_o r_o}{r_o - r_i}.
\end{equation}
At $t=0$: $\sigma_\theta \approx (0.1 \times 4.3 - 5.0 \times 4.9)/0.6 \approx -40$~MPa
(compressive under external pressure).  At $t=100$~s: $\sigma_\theta \approx -1$~MPa.

\paragraph{Results retrieval.}

\begin{lstlisting}
times      = solver.time_points()
sigh_outer = solver.clad_outer_hoop_stress()  # [n_steps] MPa
rco        = solver.clad_outer_radius()        # [n_steps] m

delta_rco_um = (rco[-1] - g.rco0) * 1e6   # outer radius change in µm
\end{lstlisting}

\begin{infobox}[Key learning point]
  Use \texttt{stress\_strain\_only} to verify the mechanical solver in
  isolation.  With temperatures fixed at $T_\text{ref}$ all thermal expansion
  strains vanish, leaving pure elastic deformation under prescribed pressure.
  This is the easiest state to compare with the Lam\'e analytical solution.
\end{infobox}

% ===========================================================================
\section{Example 3: Heat Conduction + Stress-Strain}
\label{sec:hcss}
% ===========================================================================

\textbf{File:} \texttt{examples/fred\_rod/heat\_conduction\_stress\_strain/run.py}

\medskip
Both physics blocks are active simultaneously.  The heat equation generates
a temperature distribution which drives thermal expansion, while the
mechanical equilibrium is solved at each output step using the latest
temperatures.  This is the \emph{operator-split} coupling used in FRED-ROD:
IDA advances the thermal ODE, and the mechanics are updated after each output
time.

\textbf{Scenario:}
\begin{itemize}
  \item Slow power ramp $0 \to 10^8$~W/m$^3$ over 50~s, then held constant.
  \item Coolant at 700~K, external pressure 5~MPa.
  \item Simulated to $t = 200$~s (well past steady state).
\end{itemize}

\paragraph{Why a slow power ramp?}

\begin{lstlisting}
# Ramp 0→1e8 W/m3 over 50 s, then hold
solver.set_power_density_history(
    [0.0, 50.0, 100.0, 1e6],
    [0.0, 1.0e8, 1.0e8, 1.0e8])
\end{lstlisting}

When both heat conduction and stress-strain are active, the IDA initial
condition calculation (\texttt{IDACalcIC}) must find consistent initial
derivatives for the coupled system.  A step change in power at $t=0$
introduces a large discontinuity that can cause \texttt{IDACalcIC} to fail.
Ramping the power over $\ge 50$~s gives the solver time to build a smooth
consistent initial state.

\paragraph{Enabling both physics.}

\begin{lstlisting}
solver.set_enable_heat_conduction(True)
solver.set_enable_stress_strain(True)

T0 = 700.0   # K
solver.set_initial_temperature(T0)
solver.set_coolant_temperature([0.0, 1e6], [T0, T0])
solver.set_coolant_pressure(5.0)        # MPa
solver.set_initial_gas_pressure(0.1)    # MPa
\end{lstlisting}

\paragraph{Accessing combined results.}

\begin{lstlisting}
T_arr = solver.temperatures()     # (n_steps, nf+nc)  K
sigh  = solver.clad_outer_hoop_stress()  # (n_steps,)  MPa
rco   = solver.clad_outer_radius()       # (n_steps,)  m

# Centreline and cladding inner-surface temperatures at each step:
T_ctr = T_arr[:, 0]          # innermost fuel node
T_oci = T_arr[:, g.nf]       # innermost cladding node (= clad inner surface)
\end{lstlisting}

\paragraph{Expected steady-state behaviour.}

\begin{itemize}
  \item Fuel centreline: $\approx 837$~K (44~K above outer fuel surface).
  \item Cladding outer radius grows by $\approx 17$--19~µm (thermal expansion
        from $T_0 \approx 293$~K reference to the cladding operating temperature).
  \item Hoop stress is compressive ($\approx -28$~MPa) because the external
        5~MPa coolant pressure dominates over the thermal expansion effect.
\end{itemize}

\begin{warningbox}[Power-ramp requirement]
  Always use a slow power ramp (\emph{not} a step at $t=0$) when both heat
  conduction and stress-strain are enabled.  A 50~s ramp to operating power
  is sufficient in practice.
\end{warningbox}

% ===========================================================================
\section{Example 4: Gap Close and Reopen}
\label{sec:gcr}
% ===========================================================================

\textbf{File:} \texttt{examples/fred\_rod/gapclose\_reopen/run.py}

\medskip
This example exercises the gap-state logic: the fuel and cladding can come
into physical contact (gap closes) and then separate again (gap reopens).
Root-finding events in SUNDIALS IDA detect when the gap crosses the
roughness threshold and automatically switches between the open-gap and
closed-gap residual formulations.

\textbf{Scenario (three phases):}
\begin{enumerate}
  \item \textbf{Phase 1 (0--100~s):} no power; open gap at the as-fabricated
        width of 15~µm.
  \item \textbf{Phase 2 (100--600~s):} ramp to $5 \times 10^8$~W/m$^3$,
        then hold; fuel expands, gap closes.
  \item \textbf{Phase 3 (600--1500~s):} ramp back to zero power; fuel
        contracts, gap reopens.
\end{enumerate}

The as-fabricated gap is deliberately small (15~µm instead of the 100~µm
in earlier examples) so that gap closure happens at a moderate power level.

\paragraph{Geometry with a tight as-fabricated gap.}

\begin{lstlisting}
g.rfo0 = 4.2e-3       # m
g.rci0 = 4.215e-3     # m  — 15 µm as-fabricated gap
g.ruff = 5.0e-6       # fuel roughness
g.rufc = 5.0e-6       # cladding roughness
# contact threshold = ruff + rufc = 10 µm
\end{lstlisting}

The gap is considered \emph{closed} when the mechanical gap width falls to
the combined surface roughness ($r_\text{uff} + r_\text{ufc}$).  In the
closed state the pellet-cladding contact pressure is a free variable and
the gap constraint is replaced by a no-penetration condition.

\paragraph{Multi-phase power history.}

\begin{lstlisting}
solver.set_power_density_history(
    [0.0, 100.0, 200.0, 600.0, 700.0, 1500.0],
    [0.0,   0.0, 5.0e8, 5.0e8,   0.0,    0.0])
\end{lstlisting}

\paragraph{Monitoring gap state.}

\begin{lstlisting}
gap_w = solver.gap_width()          # (n_steps,) m — axial average
pfc   = solver.contact_pressure()   # (n_steps,) MPa
rfo   = solver.fuel_outer_radius()  # (n_steps,) m
rci   = solver.clad_inner_radius()  # (n_steps,) m

# Gap-state classification
thresh = g.ruff + g.rufc    # contact threshold in m
is_closed = gap_w <= thresh  # boolean array
\end{lstlisting}

Also available in the raw IDA state vector:

\begin{lstlisting}
y_out = solver.y_out()     # (n_steps, neq)
gpres = y_out[:, 0]        # gas pressure (MPa) — first DAE variable
\end{lstlisting}

\paragraph{What to expect.}

\begin{itemize}
  \item \textbf{Phase 1:} gap $\approx 13.8$~µm (as-fabricated gap minus
        initial thermal expansion at 700~K).
  \item \textbf{Gap closure at $\approx 150$~s:} fuel expansion exceeds the
        15~µm as-fabricated gap; contact pressure jumps from 0 to $\approx 19$~MPa.
  \item \textbf{Gap reopening at $\approx 700$~s:} power removed; fuel
        contracts; gas pressure exceeds contact pressure; gap reopens.
  \item Results compare closely with the legacy FRED Fortran code
        ($|\Delta\text{gap}| < 0.2$~µm, $|\Delta\text{pfc}| < 0.1$~MPa over
        the whole 1500~s run).
\end{itemize}

\begin{infobox}[Key learning point]
  The gas pressure evolves via the ideal-gas DAE (included automatically
  in the IDA state vector).  It rises with temperature during Phase~2
  and drops when the fuel cools.  The gap reopens when the gas pressure
  exceeds the contact pressure, not simply when the fuel contracts below
  the cladding radius.
\end{infobox}

% ===========================================================================
\section{Example 5: Hot Start vs Cold Start}
\label{sec:hot-rod}
% ===========================================================================

\textbf{File:} \texttt{examples/fred\_rod/hot\_start\_demo/run.py}

\medskip
Every earlier example began at a low reference temperature and let the real
transient carry the rod up to its operating (steady-state) condition ---
either directly (Example 1) or via a manual power ramp used purely to keep
\texttt{IDACalcIC} well-behaved (Example 3).  \texttt{set\_hot\_start(True)}
replaces that workaround: the solver runs a pseudo-time pre-march at the
$t=0$ boundary conditions \emph{before} the real time loop starts, converges
to steady state, then re-anchors $t=0$ at that converged state.  The first
logged output point is then already at operating conditions.

\paragraph{Cold vs hot start.}

\begin{lstlisting}
def make_solver(hot_start: bool):
    # ... identical geometry/materials/BCs each time (solver state is not
    # resettable after run(), so build a fresh solver per call) ...
    solver = fred.FredRodSolver(g, fuel, clad, gap)
    solver.set_enable_heat_conduction(True)
    solver.set_enable_stress_strain(True)
    solver.set_hot_start(hot_start)          # False = cold start (default)
    solver.set_power_density_history([0.0, 1e6], [1.0e8, 1.0e8])  # full power at t=0
    return solver

s_cold = make_solver(hot_start=False)
s_cold.run(tend=200.0, dtout=5.0)

s_hot = make_solver(hot_start=True)
s_hot.run(tend=200.0, dtout=5.0)
\end{lstlisting}

Both runs use the identical instantaneous full-power step at $t=0$ (no
manual ramp needed --- the hot-start pseudo-time march, and IDA's own
initial-condition solve for a literal cold jump, both handle the resulting
residual imbalance without difficulty for this rod).

\paragraph{What to expect.}

\begin{itemize}
  \item \textbf{Cold run:} starts at $T_0 = 293.15$~K, rises to the
        $\approx 832.4$~K steady-state centreline temperature by
        $t \approx 15$~s.
  \item \textbf{Hot run:} already at $\approx 832.4$~K at $t=0$; the maximum
        deviation from that value over the full 200~s run is
        $\lesssim 2\times 10^{-7}$~K --- i.e.\ ``no change'', confirming the
        pre-march genuinely found the same steady state the cold run
        converges to.
\end{itemize}

\begin{warningbox}[When \emph{not} to use hot start]
  \texttt{heat\_conduction\_only} (Example 1) exists specifically to show the
  \emph{transient} heat-up over several thermal time constants --- enabling
  hot start there would erase the behaviour the example demonstrates.  More
  generally, do not use hot start for any example whose purpose is to
  observe a startup transient, a gap-closure event under power cycling
  (Example 4), or a restart/checkpoint comparison against a known reference
  trajectory (Examples 6--7): hot start changes the $t=0$ state itself, so it
  is only appropriate when steady state \emph{before} the interesting
  behaviour begins is what you want as your starting point.
\end{warningbox}

\begin{infobox}[Key learning point]
  Hot start is a numerical convenience, not new physics: FRED-ROD has no
  irradiation state to freeze during the pre-march (unlike FRED-OX/FRED-M-Na,
  Section~\ref{sec:hot-ox}, where irradiation physics is explicitly disabled
  during the march).  Use it whenever an example's own transient is not the
  point, to avoid burning simulated time (and, for stiffer coupled
  thermo-mechanical runs, solver robustness) on a ramp-up you do not care
  about.
\end{infobox}

% ===========================================================================
\section{Example 6: HDF5 Output}
\label{sec:hdf5}
% ===========================================================================

\textbf{File:} \texttt{examples/fred\_rod/hdf5\_output/run.py}

\medskip
This example demonstrates how to write results to an HDF5 file using the
\texttt{fred\_input} / \texttt{fred\_output} helper modules (located in
\texttt{fred\_platform/python/}) and how to read them back for post-processing.

\paragraph{Using \texttt{fred\_input} instead of \texttt{fred\_rod}.}

For scripts that write HDF5 output, import the validated Python wrapper:

\begin{lstlisting}
import fred_input as fred   # validated wrapper; drop-in replacement for fred_rod
import fred_output as fo    # HDF5 read/write utilities
\end{lstlisting}

\texttt{fred\_input} re-exports all the same classes (\texttt{FuelRodGeometry},
\texttt{MOX}, \texttt{T91}, \texttt{He}, \texttt{FredRodSolver}, \ldots) but
validates all inputs in Python before forwarding them to the C++ layer, giving
clearer error messages.

\paragraph{Running with HDF5 output.}

\begin{lstlisting}
solver.run(tend=1000.0, dtout=50.0, output_file='results.h5')
\end{lstlisting}

The \texttt{output\_file} argument triggers \texttt{fred\_output.write\_results\_h5()}
immediately after \texttt{run()} completes.  All results (temperatures,
mechanical state, IDA restart vectors) are written to the specified HDF5 file.

\paragraph{Two output modes.}

\begin{lstlisting}
# Mode 1: dtout mode — results at every multiple of dtout
solver.run(tend=1000.0, dtout=50.0, output_file='results.h5')

# Mode 2: all_steps mode — results at every internal IDA step
solver2.run(tend=100.0, dtout=100.0, output_file='debug.h5', all_steps=True)
\end{lstlisting}

\texttt{all\_steps=True} records every adaptive time step taken by IDA, which
can be hundreds or thousands of points per second of simulated time.  This
mode is useful for diagnosing solver behaviour but produces large files.

\paragraph{HDF5 file layout.}

The file written by \texttt{write\_results\_h5()} has the following structure:

\begin{lstlisting}[language=bash]
results.h5
  /metadata          -- attributes: app, nz, nf, nc, n_steps, all_steps
  /time              -- [n_steps]                    seconds
  /thermal/
      T              -- [n_steps, nz, nf+nc]         K
      peak_T_fuel    -- [n_steps]                    K
      gap_width      -- [n_steps]                    m  (axial avg)
  /mechanical/
      pfc            -- [n_steps]                    MPa
      rfo            -- [n_steps]                    m
      rci            -- [n_steps]                    m
      rco            -- [n_steps]                    m
      sigh_outer     -- [n_steps]                    MPa
  /restart/
      y              -- [n_steps, neq]               IDA state
      yp             -- [n_steps, neq]               IDA derivatives
\end{lstlisting}

\paragraph{Reading results back.}

\begin{lstlisting}
r = fo.read_results_h5('results.h5')

# Metadata
nz, nf, nc = r['nz'], r['nf'], r['nc']

# Arrays
times  = r['time']           # [n_steps] s
T      = r['T']              # [n_steps, nz, nf+nc] K
peak_T = r['peak_T_fuel']    # [n_steps] K
gap_w  = r['gap_width']      # [n_steps] m
pfc    = r['pfc']            # [n_steps] MPa
sigh   = r['sigh_outer']     # [n_steps] MPa

# Access radial temperature profile at final time, first axial layer:
T_profile = T[-1, 0, :]      # shape (nf+nc,)
T_fuel    = T[-1, 0, :nf]    # fuel nodes
T_clad    = T[-1, 0, nf:]    # cladding nodes

# IDA restart vectors (used by load_checkpoint / load_snapshot):
y_last  = r['restart_y'][-1, :]   # y-vector at final step
yp_last = r['restart_yp'][-1, :]
\end{lstlisting}

\begin{infobox}[Key learning point]
  Storing the IDA \texttt{y}/\texttt{yp} vectors in the HDF5 file enables
  post-run analysis: you can reconstruct the solver at any saved output time
  and inspect or restart from that state.  The \texttt{restart\_snapshot}
  example (Section~\ref{sec:rs}) shows a more structured API for this.
\end{infobox}

% ===========================================================================
\section{Example 7: Restart and Snapshot}
\label{sec:rs}
% ===========================================================================

\textbf{File:} \texttt{examples/fred\_rod/restart\_snapshot/}

\medskip
This example contains three scripts that illustrate two complementary
persistence mechanisms:

\begin{description}
  \item[Checkpoint] A single overwriting file for fault recovery.  Simulation
        time is preserved on reload.
  \item[Snapshot] A permanent restart-state file.  Physical state (temperatures,
        strains, gap) is preserved but simulation time is reset to~0.  Used to
        start a new transient from a pre-conditioned fuel state.
\end{description}

\paragraph{Script 0: Full reference run.}

\texttt{0\_full\_run.py} runs the complete simulation from $t=0$ to $t=5000$~s
with both checkpointing and snapshot saving enabled.

\begin{lstlisting}
WORK_DIR = 'ckpt'
tend   = 5000.0  # s
t_half = 3000.0  # s — mid-point snapshot

s = make_solver()

# Enable fault-recovery checkpoint (overwritten every dtout)
s.set_checkpoint_prefix(os.path.join(WORK_DIR, 'rod'))

# Save a snapshot at t_half; the final snapshot is always saved automatically
s.set_snapshot_prefix(os.path.join(WORK_DIR, 'rod'),
                      snapshot_times=[t_half])

s.run(tend, dtout=1000.0)
\end{lstlisting}

After running, the \texttt{ckpt/} directory contains:
\begin{lstlisting}[language=bash]
rod_checkpoint.chk     # fault-recovery checkpoint (overwritten at each dtout)
rod_frame1.snapshot    # snapshot saved at t_half = 3000 s
rod_frame2.snapshot    # snapshot saved at tend  = 5000 s  (automatic)
ref_times.npy          # reference output times (for comparison in scripts 1 & 2)
ref_T.npy              # reference peak fuel temperatures
ref_gap.npy            # reference gap widths
\end{lstlisting}

\paragraph{Script 1: Checkpoint restart.}

\texttt{1\_checkpoint\_restart.py} demonstrates fault recovery.  It simulates a
crash at \texttt{t\_half} by running only to that point, then reloads the
checkpoint and continues to \texttt{tend}.

\begin{lstlisting}
# Step A: partial run 0 → t_half (crashes here)
s_partial = make_solver()
s_partial.set_checkpoint_prefix(ckpt_prefix)
s_partial.run(t_half, dtout)

# Step B: reload checkpoint and continue to tend
s_restart = make_solver()
s_restart.set_checkpoint_prefix(ckpt_prefix)
s_restart.load_checkpoint(ckpt_file)
print(f"Restart time: {s_restart.restart_time():.0f} s")

s_restart.run(tend, dtout)
\end{lstlisting}

\texttt{restart\_time()} returns the time from which the simulation will
continue ($= t_\text{half}$ in this case).  Results from the restarted run
match the reference run to within $10^{-6}$ relative error on all quantities.

\paragraph{Script 2: Continuation via snapshot.}

\texttt{2\_continuation.py} demonstrates the physical-state snapshot for a
two-phase simulation: an ``irradiation'' phase followed by a ``transient''
phase with fresh time coordinates.

\begin{lstlisting}
# Phase 1: irradiation, 0 → t_half
s1 = make_solver()
s1.set_snapshot_prefix(snap_prefix)   # saves rod_frame1.snapshot at end
s1.run(t_half, dtout)

# Phase 2: fresh start from the irradiated state
s2 = make_solver(t_end=t_phase2)
s2.load_snapshot(snap_file)
print(f"restart_time = {s2.restart_time():.1f} s")  # prints 0.0 — time reset!

s2.run(t_phase2, dtout)

# Stitch results: Phase 2 times are offset by t_half for plotting
t_combined = np.concatenate([t1, t2[1:] + t_half])
\end{lstlisting}

The key difference between checkpoint and snapshot is the time-reset:
\begin{itemize}
  \item \texttt{load\_checkpoint()} $\to$ \texttt{restart\_time()} returns
        the saved physical time; the run continues from there.
  \item \texttt{load\_snapshot()} $\to$ \texttt{restart\_time()} returns
        $0.0$; the run restarts from $t=0$ with the pre-irradiated fuel state.
\end{itemize}

\begin{infobox}[Key learning point]
  Snapshots are the FRED-ROD mechanism for multi-stage simulations: you can
  build up a library of pre-conditioned fuel states (different burnup levels,
  gap sizes, strain states) and quickly start transient calculations from
  each one without re-running the irradiation phase.
\end{infobox}

% ===========================================================================
\section{Example 8: Python-Defined Custom Materials}
\label{sec:pymat}
% ===========================================================================

\textbf{File:} \texttt{examples/fred\_rod/python\_materials/run.py}

\medskip
FRED-ROD exposes its abstract material interfaces to Python via
\texttt{pybind11} trampoline classes.  Any class that inherits from
\texttt{fred.FuelPelletMaterial} or \texttt{fred.CladdingMaterial} and
overrides the required methods can be passed directly to
\texttt{FredRodSolver}.  No C++ recompilation is needed.

This is the recommended approach for rapid prototyping of new material
correlations before promoting them to a permanent C++ implementation.

\subsection{Subclassing \texttt{FuelPelletMaterial}}

\begin{lstlisting}
import fred_rod as fred

T_REF = 293.15   # K — reference temperature for thermal expansion

class GenericOxideFuel(fred.FuelPelletMaterial):
    """
    Simplified oxide fuel with analytic correlations.
    All units follow the FRED convention (K, MPa, kg/m3, W/mK, J/kgK).
    """

    def thermalConductivity(self, T):
        """k(T) = 1/(A + B*T) + C*T^3  [W/(m·K)]"""
        A, B, C = 0.045, 2.8e-4, 5.6e-11
        return 1.0 / (A + B * T) + C * T**3

    def heatCapacity(self, T):
        """cp(T) = 260 + 0.10*T  [J/(kg·K)]"""
        return 260.0 + 0.10 * T

    def thermalExpansionStrain(self, T):
        """Linear CTE = 1e-5 /K"""
        return 1.0e-5 * (T - T_REF)

    def youngsModulus(self, T, density):
        """E0 = 200 GPa, linear softening  [MPa]"""
        return max(1.0e4, 2.0e5 * (1.0 - 1.5e-4 * (T - T_REF)))

    def poissonRatio(self):
        return 0.30

    def referenceDensity(self):
        return 10500.0    # kg/m³  (as-fabricated, ~95.8% TD)

    def theoreticalDensity(self):
        return 10960.0    # kg/m³

    def meltingTemperature(self):
        return 3100.0     # K
\end{lstlisting}

The \texttt{thermalConductivity(T)} form
$k = 1/(A + BT) + CT^3$ is a standard two-term expression for oxide ceramics:
the first term captures phonon scattering (decreasing with temperature) and
the second term captures the small-polaron contribution (increasing with
temperature).

\subsection{Subclassing \texttt{CladdingMaterial}}

\begin{lstlisting}
class GenericSteelClad(fred.CladdingMaterial):

    def thermalConductivity(self, T):
        """k(T) = 23.5 + 0.005*(T - 273.15)  [W/(m·K)]"""
        return 23.5 + 0.005 * (T - 273.15)

    def heatCapacity(self, T):
        """cp(T) = 480 + 0.12*T  [J/(kg·K)]"""
        return 480.0 + 0.12 * T

    def thermalExpansionStrain(self, T):
        """CTE = 11.5e-6 /K"""
        return 11.5e-6 * (T - T_REF)

    def youngsModulus(self, T):
        """Linear softening from 200 GPa  [MPa]"""
        return max(1.0e4, 2.0e5 - 60.0 * (T - T_REF))

    def poissonRatio(self):
        return 0.30

    def meyerHardness(self, T):
        """H_M(T) = 5e9 * T^(-0.2)  [Pa]"""
        return 5.0e9 * T**(-0.2)

    def referenceDensity(self):
        return 7800.0   # kg/m³
\end{lstlisting}

Note that \texttt{CladdingMaterial.youngsModulus(T)} takes only temperature
(no density), whereas \texttt{FuelPelletMaterial.youngsModulus(T, density)}
takes both (porosity affects the fuel modulus).  \texttt{meyerHardness(T)}
is used in the Ross--Stoute solid contact conductance model when the gap is
closed.

\subsection{Required abstract methods}

Table~\ref{tab:abstract} lists every method that \emph{must} be overridden.
FRED will raise a \texttt{RuntimeError} if any method is not implemented.

\begin{table}[h]
\centering
\caption{Abstract methods that must be implemented in Python subclasses.}
\label{tab:abstract}
\begin{tabular}{lll}
\toprule
Class & Method & Returns \\
\midrule
\multirow{8}{*}{\texttt{FuelPelletMaterial}}
  & \texttt{thermalConductivity(T)}       & W/(m$\cdot$K) \\
  & \texttt{heatCapacity(T)}              & J/(kg$\cdot$K) \\
  & \texttt{thermalExpansionStrain(T)}    & --- (strain, zero at $T_\text{ref}$) \\
  & \texttt{youngsModulus(T, density)}    & MPa \\
  & \texttt{poissonRatio()}               & --- \\
  & \texttt{referenceDensity()}           & kg/m$^3$ \\
  & \texttt{theoreticalDensity()}         & kg/m$^3$ \\
  & \texttt{meltingTemperature()}         & K \\
\midrule
\multirow{7}{*}{\texttt{CladdingMaterial}}
  & \texttt{thermalConductivity(T)}       & W/(m$\cdot$K) \\
  & \texttt{heatCapacity(T)}             & J/(kg$\cdot$K) \\
  & \texttt{thermalExpansionStrain(T)}   & --- \\
  & \texttt{youngsModulus(T)}            & MPa \\
  & \texttt{poissonRatio()}              & --- \\
  & \texttt{meyerHardness(T)}           & Pa \\
  & \texttt{referenceDensity()}         & kg/m$^3$ \\
\bottomrule
\end{tabular}
\end{table}

\subsection{Using the custom materials in a simulation}

\begin{lstlisting}
fuel = GenericOxideFuel()
clad = GenericSteelClad()
gap  = fred.DummyGapMaterial()

solver = fred.FredRodSolver(g, fuel, clad, gap)  # standard API
solver.set_enable_heat_conduction(True)
solver.set_enable_stress_strain(True)

# ... set BCs and run as usual ...
solver.run(tend=1000.0, dtout=100.0)
\end{lstlisting}

There is no difference in how \texttt{FredRodSolver} is used: the Python
objects are polymorphically identical to the built-in C++ materials from the
solver's perspective.

\subsection{Verifying your material with the built-in API}

You can call the material methods directly from Python to verify they are
correct before running a full simulation:

\begin{lstlisting}
fuel = GenericOxideFuel()

for T in [400, 700, 1000, 1400]:
    k  = fuel.thermal_conductivity(T)
    cp = fuel.heat_capacity(T)
    print(f"T={T:4d} K: k={k:.3f} W/mK, cp={cp:.1f} J/kgK")
\end{lstlisting}

\begin{warningbox}[Python call overhead]
  Each Python material method is called by the C++ solver once per radial
  node per Newton iteration per time step.  For a run with thousands of time
  steps and fine meshes this can introduce measurable overhead compared to a
  compiled C++ material.  Once a correlation is validated, implement it in
  C++ and add pybind11 bindings following the pattern in
  \texttt{src/apps/fred\_rod/fuelpelletmaterial/UO2.hpp/.cpp} and
  \texttt{bindings.cpp}.
\end{warningbox}

% ===========================================================================
\section{Example 9: Advanced Prototyping --- Python to C++}
\label{sec:proto}
% ===========================================================================

\textbf{Files:} \texttt{examples/fred\_rod/advanced\_prototyping/}

\medskip
Section~\ref{sec:pymat} showed how to prototype fuel and cladding materials in
Python.  The same mechanism applies to gap materials via \texttt{fred.GapMaterial}.
This example uses a physically motivated model --- a burnup-dependent He/Kr/Xe
gap conductivity --- to demonstrate a two-step development loop:

\begin{enumerate}
  \item \textbf{Step 1} (\texttt{run\_python.py}): implement the model as a
        Python subclass.  No rebuild required; changes take effect immediately.
  \item \textbf{Step 2} (\texttt{run\_cpp.py}): use the compiled C++ equivalent
        \texttt{fred.HeKrXe}.  Results are numerically identical; C++ removes
        Python call overhead in the inner solver loop.
\end{enumerate}

\subsection{Physics: fission gas release and gap conductivity}

Fresh fuel pins are filled with helium (excellent thermal conductor,
$k_\text{He} \approx 0.27$~W/mK at 700~K).  During irradiation, xenon and
krypton --- fission products with conductivities 10--20$\times$ lower than
helium --- are released into the gap gas.  Even a moderate fission gas
release (FGR) fraction significantly raises the fuel temperature.

The model implemented here uses three steps:

\paragraph{FGR fraction (Waltar-Reynolds, restructured zone):}
\begin{align}
  B_\text{MWd} &= B_\text{at\%} \times 9.4
  \qquad (1\,\text{at\%} \approx 9.4\,\text{MWd/kgHM}) \notag\\
  \text{FGR} &= \max\!\left(0,\; 1 - \frac{4.7}{B_\text{MWd}}
                \left(1 - e^{-B_\text{MWd}/5.9}\right)\right)
\end{align}

\paragraph{Molar balance:}
\begin{align}
  n_{\text{He},0} &= \frac{P_\text{fill} V_\text{free}}{RT_\text{ref}}
  \quad\approx 6.15\times10^{-5}\,\text{mol}
  \quad (P=0.1\,\text{MPa},\; V=1.5\,\text{cm}^3,\; T=293\,\text{K})
  \notag\\
  n_\text{FG} &= \text{FGR} \times 0.25 \times
                \frac{B_\text{at\%}}{100} \times \frac{m_\text{fuel}}{M_\text{HM}}
  \notag\\
  y_\text{He} &= (n_{\text{He},0} + 0.0385\,n_\text{FG})\;/\;n_\text{total}
  \notag\\
  y_\text{Kr} &= 0.0769\,n_\text{FG}\;/\;n_\text{total},
  \qquad
  y_\text{Xe} = 0.8846\,n_\text{FG}\;/\;n_\text{total}
\end{align}

\paragraph{Geometric-mean mixture conductivity (FRED.f90 \texttt{gaphtc}):}
\begin{equation}
  k_\text{mix}(T) = k_\text{He}^{y_\text{He}}\;
                    k_\text{Kr}^{y_\text{Kr}}\;
                    k_\text{Xe}^{y_\text{Xe}},
  \qquad k_i(T) = A_i\, T^{B_i}
\end{equation}

\noindent Typical values at 700~K and several burnup levels:

\begin{center}
\begin{tabular}{rrrrrr}
\toprule
$B$ [at\%] & FGR [\%] & $y_\text{He}$ & $y_\text{Kr}$ & $y_\text{Xe}$
           & $k_\text{mix}$ [W/mK] \\
\midrule
0.0 & 0   & 1.000 & 0.000 & 0.000 & 0.274 \\
1.0 & 60  & 0.513 & 0.039 & 0.448 & 0.062 \\
5.0 & 90  & 0.149 & 0.068 & 0.783 & 0.020 \\
\bottomrule
\end{tabular}
\end{center}

\subsection{Step 1: Python prototype (\texttt{run\_python.py})}

The \texttt{GapMaterial} interface requires a single method:

\begin{lstlisting}
import fred_rod as fred
import math

class HeKrXePy(fred.GapMaterial):

    _A_He, _B_He = 2.639e-3, 0.7085   # k_He(T) = A * T^B  [W/(m·K)]
    _A_Kr, _B_Kr = 8.247e-5, 0.8363
    _A_Xe, _B_Xe = 4.351e-5, 0.8618

    def __init__(self, bup_atpct=0.0):
        super().__init__()
        bup_MWd = bup_atpct * 9.4
        fgr = 0.0
        if bup_MWd > 0.0:
            a   = 4.7 / bup_MWd * (1.0 - math.exp(-bup_MWd / 5.9))
            fgr = max(0.0, 1.0 - a)

        n_He0  = 0.1e6 * 1.5e-6 / (8.314 * 293.15)   # ~6.15e-5 mol
        fggen  = 0.25 * (bup_atpct / 100.0) * 0.010 / 0.238
        fgrel  = fgr * fggen
        n_Xe   = 0.8846 * fgrel
        n_Kr   = 0.0769 * fgrel
        n_HeFG = 0.0385 * fgrel
        n_tot  = n_He0 + n_HeFG + n_Kr + n_Xe

        self.y_He = (n_He0 + n_HeFG) / n_tot
        self.y_Kr = n_Kr / n_tot
        self.y_Xe = n_Xe / n_tot

    def gapConductivity(self, T):
        k_He = self._A_He * T ** self._B_He
        k_Kr = self._A_Kr * T ** self._B_Kr
        k_Xe = self._A_Xe * T ** self._B_Xe
        return k_He ** self.y_He * k_Kr ** self.y_Kr * k_Xe ** self.y_Xe
\end{lstlisting}

The platform takes \texttt{gapConductivity(T)} and converts it to a gap
conductance internally using the Lanning-Hann gas-jump model with the
actual (time-varying) gap width and gas pressure.  The material class does
not need to know the gap width.

\begin{infobox}[\texttt{GapMaterial} interface]
  Only \texttt{gapConductivity(T)} must be overridden.  The solver calls it
  at each Newton iteration to obtain $k$ [W/(m$\cdot$K)], then computes the
  full gap conductance $h$ [W/(m$^2\cdot$K)] using the gap geometry.
  Optionally override \texttt{isGasBond()} (\texttt{True} by default) and
  \texttt{clampConductance(h, gap\_closed)} for liquid-bond materials
  (sodium, lead) where the conversion from $k$ to $h$ is different.
\end{infobox}

Three simulations are run with this gap material at increasing burnup levels,
with UO2 fuel, AIM1 cladding, and a constant volumetric power of
$3\times10^8$~W/m$^3$:

\begin{center}
\begin{tabular}{lrrr}
\toprule
Burnup & $T_\text{peak}$ [K] & Gap [$\mu$m] & $\sigma_{\theta,\text{outer}}$ [MPa] \\
\midrule
0.0 at\% (as-fabricated) & 1272 & 81.8 & 22.3 \\
1.0 at\% (moderate)      & 1678 & 63.4 & 22.3 \\
5.0 at\% (high)          & 2008 & 46.1 & 22.4 \\
\bottomrule
\end{tabular}
\end{center}

The $+735$~K rise from 0 to 5~at\% illustrates the strong coupling between
fission gas release and peak fuel temperature.

\subsection{Step 2: Compiled C++ equivalent (\texttt{run\_cpp.py})}

Once the Python model is validated, use the compiled \texttt{fred.HeKrXe}
class --- which implements the same physics in C++:

\begin{lstlisting}
gap = fred.HeKrXe(bup_atpct=5.0)   # compiled class, same physics
s   = fred.FredRodSolver(geom, fred.UO2(), fred.AIM1(), gap)
# ... identical setup and run() call ...
\end{lstlisting}

The \texttt{run\_cpp.py} script also verifies that the conductivity values
agree with the Python prototype to machine precision before running the
simulations.

\subsection{Adding your own compiled gap material}

To follow the same Python$\to$C++ path for your own model:

\begin{enumerate}
  \item Write a header-only class in
        \texttt{src/apps/fred\_rod/gapmaterial/MyGap.hpp}
        that inherits \texttt{fred::GapMaterial} and overrides
        \texttt{gapConductivity(double T) const}.
  \item Add a pybind11 binding block in
        \texttt{src/apps/fred\_rod/bindings.cpp}:
        \begin{lstlisting}
#include "apps/fred_rod/gapmaterial/MyGap.hpp"
// inside bind_fred_rod():
py::class_<MyGap, GapMaterial>(m, "MyGap")
    .def(py::init<double>(), py::arg("param") = 0.0,
         "Description of MyGap.");
        \end{lstlisting}
  \item Rebuild: \texttt{make fred-rod}.
        No \texttt{Makefile} source-list change is needed for header-only classes.
\end{enumerate}

See \texttt{examples/fred\_rod/advanced\_prototyping/README.md} for the
complete walkthrough, including the \texttt{.cpp}-file variant and liquid-bond
material instructions.

% ===========================================================================
\section{FRED-ROD API Reference}
\label{sec:api-rod}
% ===========================================================================

\subsection{Geometry: \texttt{FuelRodGeometry}}

\begin{center}
\begin{tabular}{lll}
\toprule
Attribute & Type & Description \\
\midrule
\texttt{nf}   & int & Fuel radial nodes ($\ge 2$) \\
\texttt{nc}   & int & Cladding radial nodes ($\ge 2$) \\
\texttt{nz}   & int & Axial layers ($\ge 1$) \\
\texttt{rfi0} & float & Inner fuel radius [m] (0 = solid pellet) \\
\texttt{rfo0} & float & Outer fuel radius [m] \\
\texttt{rci0} & float & Inner cladding radius [m] \\
\texttt{rco0} & float & Outer cladding radius [m] \\
\texttt{dz0}  & list  & Axial layer heights [m] (length \texttt{nz}) \\
\texttt{vgp}  & float & Gas plenum volume [m$^3$] \\
\texttt{ruff} & float & Fuel outer surface roughness [m] \\
\texttt{rufc} & float & Cladding inner surface roughness [m] \\
\bottomrule
\end{tabular}
\end{center}

\subsection{Built-in materials}

\begin{center}
\begin{tabular}{lll}
\toprule
Class & Type & Notes \\
\midrule
\texttt{UO2(reference\_density)}         & Fuel    & Philipponneau conductivity \\
\texttt{MOX(pu\_content, [ref\_density])} & Fuel    & Philipponneau; $x_\text{Pu} \in [0.01,0.5]$ \\
\texttt{DummyFuelPellet()}               & Fuel    & Constant $k=10$, $\rho=10000$, $c_p=100$ \\
\texttt{AIM1([reference\_density])}       & Cladding & Austenitic SS (Luzzi correlations) \\
\texttt{T91([reference\_density])}        & Cladding & Ferritic-martensitic (KIT) \\
\texttt{DummyCladding()}                  & Cladding & Constant $k=10$, $\rho=10000$, $c_p=100$ \\
\texttt{He()}                             & Gap      & Pure helium \\
\texttt{HeKrXe(bup\_atpct)}               & Gap      & He-Kr-Xe mixture; composition from burnup \\
\texttt{DummyGapMaterial()}               & Gap      & $h(T) = 500 + 2.5T$ [W/m$^2$K] \\
\bottomrule
\end{tabular}
\end{center}

\subsection{Solver methods: \texttt{FredRodSolver}}

\begin{center}
\begin{tabular}{ll}
\toprule
Method & Description \\
\midrule
\texttt{set\_enable\_heat\_conduction(bool)}       & Toggle thermal ODE \\
\texttt{set\_enable\_stress\_strain(bool)}         & Toggle mechanics \\
\texttt{set\_initial\_temperature(T0\_K)}          & Uniform IC [K] \\
\texttt{set\_initial\_gas\_pressure(p0\_MPa)}      & Fill-gas pressure IC [MPa] \\
\texttt{set\_coolant\_temperature(times, T\_K)}    & Outer BC history [K] \\
\texttt{set\_coolant\_pressure(p\_MPa)}            & Constant coolant pressure [MPa] \\
\texttt{set\_coolant\_pressure\_history(t, p)}     & Time-varying coolant pressure [MPa] \\
\texttt{set\_power\_density\_history(times, qv)}   & Volumetric heat [W/m$^3$] \\
\texttt{set\_tolerances(rtol, atol)}               & IDA tolerances \\
\texttt{set\_hot\_start(bool)}                     & Pseudo-time pre-march to steady state before $t=0$ (Section~\ref{sec:hot-rod}) \\
\texttt{set\_checkpoint\_prefix(prefix)}           & Enable fault-recovery checkpoint \\
\texttt{set\_snapshot\_prefix(prefix, [times])}    & Enable physical-state snapshots \\
\texttt{load\_checkpoint(filename)}                & Load checkpoint (time preserved) \\
\texttt{load\_snapshot(filename)}                  & Load snapshot (time reset to 0) \\
\texttt{set\_output\_file(filename)}               & Stream results to HDF5 during \texttt{run()} (native path) \\
\texttt{run(tend, dtout, [all\_steps])}            & Execute simulation \\
\midrule
\texttt{time\_points()}               & Output times [s] \\
\texttt{temperatures()}               & Nodal temperatures, shape \texttt{(n,nz*(nf+nc))} [K] \\
\texttt{peak\_fuel\_temperature()}    & Peak fuel temperature [K] \\
\texttt{gap\_width()}                 & Axial-avg gap width [m] \\
\texttt{contact\_pressure()}          & Axial-avg contact pressure [MPa] \\
\texttt{clad\_outer\_hoop\_stress()}  & Axial-avg outer clad hoop stress [MPa] \\
\texttt{clad\_outer\_radius()}        & Axial-avg outer cladding radius [m] \\
\texttt{fuel\_outer\_radius()}        & Axial-avg fuel outer radius [m] \\
\texttt{clad\_inner\_radius()}        & Axial-avg cladding inner radius [m] \\
\texttt{y\_out()}                     & IDA state vector, shape \texttt{(n,neq)} \\
\texttt{yp\_out()}                    & IDA derivative vector, shape \texttt{(n,neq)} \\
\texttt{restart\_time()}              & Restart time from last load [s] \\
\bottomrule
\end{tabular}
\end{center}

% ===========================================================================
\section{FRED-ROD Troubleshooting}
\label{sec:trouble-rod}
% ===========================================================================

\begin{description}[leftmargin=0pt]
  \item[\texttt{IDACalcIC} failure at $t=0$ with heat+stress-strain enabled.]
    Either use a slow power ramp (at least 50~s) instead of a step at $t=0$
    (see Section~\ref{sec:hcss}), or --- if a startup transient is not what
    you want to observe anyway --- switch to \texttt{set\_hot\_start(True)}
    (Section~\ref{sec:hot-rod}), which sidesteps the issue entirely by
    pre-converging before $t=0$.

  \item[Gap does not close when expected.]
    Check that \texttt{rci0 - rfo0} is small enough relative to the
    operating temperature difference, and that \texttt{ruff + rufc} (the
    contact threshold) is less than the as-fabricated gap.

  \item[Custom Python material causes segfault.]
    Ensure your subclass does not raise exceptions inside the overridden
    methods: C++ does not handle Python exceptions gracefully.  Validate
    the correlation range in Python before passing the material to the solver.

  \item[HDF5 file not created.]
    The \texttt{output\_file} argument to \texttt{run()} requires
    \texttt{h5py} to be installed (\texttt{pip install h5py}).

  \item[\texttt{fred\_input} or \texttt{fred\_output} not found.]
    These modules live in \texttt{fred\_platform/python/}.  Either add that
    directory to your \texttt{PYTHONPATH}, copy the files to your working
    directory, or add \texttt{sys.path.insert(0, '/path/to/fred\_platform/python')}
    to your script.
\end{description}

% ===========================================================================
\part{FRED-OX}
\label{part:ox}
% ===========================================================================

FRED-OX uses the same coupled heat-conduction / thermo-elastic stress-strain
equations as FRED-ROD (Part~\ref{part:rod}) and adds oxide-fuel-specific
irradiation physics on top: empirical base-irradiation correlations for
fuel/cladding properties, and fission-gas production/release, which feed
back into gas-plenum pressure (an extra DAE state) and gap conductance.
Where FRED-ROD leaves irradiation modelling entirely to the user, FRED-OX
provides it built in for oxide (UO$_2$/MOX) fuel.

The solver class is \texttt{FredOxSolver} (module \texttt{fred\_ox} /
compiled extension \texttt{\_fred\_ox}), used in place of
\texttt{FredRodSolver}, with FRED-OX-specific material classes:
\texttt{FredOxMOX(pu\_content, rof0, sto0)}, \texttt{FredOxT91(reference\_density)},
\texttt{FredOxGapMaterial()}.  Geometry is still built with the same
\texttt{FuelRodGeometry}.

\section{Recommended Example Order}
\label{sec:ox-order}

The example scripts under \texttt{examples/fred\_ox/} are best worked
through in the following order:

\begin{enumerate}
  \item \textbf{\texttt{base\_irradiation}} --- first FRED-OX run: thermal +
        base-irradiation + fission-gas-release physics, validated against
        the legacy FRED Fortran code.
  \item \textbf{\texttt{hot\_start\_demo}} --- cold vs hot start for FRED-OX,
        with irradiation physics explicitly frozen during the pre-march.
  \item \textbf{\texttt{hdf5\_output}} --- the FRED-OX-specific HDF5 fields
        (gas pressure, fission gas, burnup) on top of temperature.
  \item \textbf{\texttt{restart}} --- checkpoint-based fault recovery,
        carrying irradiation state (burnup, gas inventory) through a restart.
  \item \textbf{\texttt{restart\_snapshot}} --- checkpoint \emph{and}
        named-snapshot workflows combined.
  \item \textbf{\texttt{ESFR\_pin}} --- the most advanced example: a 600-day
        base-irradiation benchmark followed by two fast transients (ULOF,
        UTOP) loaded from a saved irradiated snapshot.
\end{enumerate}

\begin{infobox}[Why this order?]
  \texttt{base\_irradiation} is the FRED-OX analogue of FRED-ROD's Examples
  1--4: it is the first place you see the extra irradiation state (burnup,
  fission gas, gas pressure) and the per-layer (axially resolved) boundary
  conditions that most real FRED-OX runs use.  \texttt{hot\_start\_demo}
  builds directly on it (same style of pin, shorter run) and is easiest to
  understand once you know what ``irradiation effects'' means for FRED-OX.
  \texttt{hdf5\_output}, \texttt{restart}, and \texttt{restart\_snapshot}
  mirror the equivalent FRED-ROD examples (Sections~\ref{sec:hdf5},
  \ref{sec:rs}) but are worth revisiting for OX specifically because
  irradiation state (burnup, gas inventory) must persist correctly across
  restarts, not just temperatures and stresses.  \texttt{ESFR\_pin} is last
  because it combines everything: a long base-irradiation run, snapshotting,
  and transient analysis starting from a realistically irradiated state.
\end{infobox}

% ===========================================================================
\section{Example OX-1: Base Irradiation}
\label{sec:ox-base}
% ===========================================================================

\textbf{File:} \texttt{examples/fred\_ox/base\_irradiation/irradiate\_10d/run.py}

\medskip
This is the FRED-OX equivalent of FRED-ROD's Example 3
(Section~\ref{sec:hcss}): heat conduction and stress-strain are both active,
but now base-irradiation correlations and fission-gas release are active
too, and boundary conditions are specified \emph{per axial layer} rather
than as single rod-wide histories.

\textbf{Scenario:} an ESFR-SMART-style pin, \texttt{nz=20} axial layers
(\texttt{dz=0.05~m} each, 1.0~m active length), \texttt{nf=22} fuel nodes,
\texttt{nc=5} cladding nodes.  MOX fuel (\texttt{pu\_content=0.1799},
\texttt{rof0=10542~kg/m$^3$}), T91 cladding
(\texttt{reference\_density=7790~kg/m$^3$}); \texttt{rfi0=1.56e-3~m},
\texttt{rfo0=4.68e-3~m}, \texttt{rci0=rfo0+1.55e-4~m}, \texttt{rco0=5.335e-3~m},
\texttt{vgp=7.19e-5~m$^3$}.  Power ramps from 0 to a per-layer peak of
$6.934\times10^8$~W/m$^3$ (layer 12) between $t=3600$~s and $t=86400$~s;
clad outer temperature ramps from 668.0~K up to 824.31~K over the same
window; \texttt{coolant\_pressure=0.2~MPa}, \texttt{initial\_gas\_pressure=0.1~MPa};
\texttt{tend=86400~s}, \texttt{dtout=1800~s}; tolerances \texttt{(1e-5, 1e-7)}.

\paragraph{Per-layer boundary conditions.}

\begin{lstlisting}
solver.set_power_density_history_per_layer(times, qqv_per_layer)   # W/m3, one column per layer
solver.set_coolant_temperature_per_layer(times, Tco_per_layer)     # K
solver.set_plenum_temperature_history(times, T_plenum)             # K
solver.set_swelling_multiplier(0.8)   # scales solid fission-product swelling
\end{lstlisting}

\texttt{set\_power\_density\_history\_per\_layer} and
\texttt{set\_coolant\_temperature\_per\_layer} take a 2-D array (time points
$\times$ \texttt{nz}) instead of the single-column history used in FRED-ROD;
this is what lets the axial power/temperature profile vary along the pin.
\texttt{set\_plenum\_temperature\_history} is new: it feeds the gas-plenum
temperature used in the gas-pressure DAE.

\paragraph{New result accessors.}

\begin{lstlisting}
gpres = solver.gas_pressure()          # (n_steps,) MPa  -- plenum gas pressure
fggen = solver.fg_generated()          # (n_steps,) mol  -- total fission gas generated
fgrel = solver.fg_released()           # (n_steps,) mol  -- total fission gas released
bup   = solver.burnup()                # (n_steps,) MWd/kgU
gap_w = solver.gap_width()             # (n_steps,) m, axial average
Tpeak = solver.peak_fuel_temperature() # (n_steps,) K
\end{lstlisting}

\paragraph{Legacy comparison at $t=86400$~s.}

\begin{center}
\begin{tabular}{lrr}
\toprule
Quantity & FRED-OX & Legacy FRED (Fortran) \\
\midrule
Gas pressure [MPa]                 & 0.24326 & --- \\
Fission gas generated [cm$^3$ STP] & 0.33828 & --- \\
Fission gas released [cm$^3$ STP]  & 0.013437 & --- \\
FGR fraction [\%]                  & 3.972 & (reference: \texttt{outfrd000000000048}) \\
\bottomrule
\end{tabular}
\end{center}

Gas amounts are stored in moles internally and converted to cm$^3$ at STP for
comparison with the legacy output via
$V_\text{STP} = n \times 8.314 \times 293.15 / 10^5 \times 10^6$.  The script
produces \texttt{plots/gas\_pressure.png}, \texttt{plots/peak\_T\_fuel.png},
\texttt{plots/fission\_gas.png}, and \texttt{plots/gap\_width.png}.

\begin{infobox}[Key learning point]
  \texttt{set\_swelling\_multiplier} scales only the \emph{solid}
  fission-product swelling correlation --- a tuning knob for matching a
  specific fuel design's densification/swelling behaviour.  It does not
  affect gas-bubble swelling or fission-gas release, which are computed from
  the burnup history directly.
\end{infobox}

% ===========================================================================
\section{Example OX-2: Hot Start vs Cold Start}
\label{sec:hot-ox}
% ===========================================================================

\textbf{File:} \texttt{examples/fred\_ox/hot\_start\_demo/run.py}

\medskip
The same \texttt{set\_hot\_start(True)} mechanism introduced for FRED-ROD
(Section~\ref{sec:hot-rod}), but here the pseudo-time pre-march explicitly
\textbf{disables irradiation physics} (fission gas release, swelling,
burnup-dependent conductivity) while it converges the thermal/mechanical
state, then re-enables irradiation and re-anchors $t=0$.  This matters
because FRED-OX's irradiation state is time-integrated (burnup, gas
inventory) --- letting it accumulate during an artificial pseudo-time march
would give a physically meaningless burnup at ``$t=0$''.

\textbf{Scenario:} a smaller single-layer MOX pin (\texttt{nz=1, nf=10,
nc=3}) so the demo runs fast: \texttt{rfo0=4.68e-3~m},
\texttt{rci0=rfo0+1.55e-4~m} (155~$\mu$m gap), \texttt{rco0=5.335e-3~m},
\texttt{vgp=7.19e-5~m$^3$}; power $2\times10^8$~W/m$^3$; coolant and plenum
temperature both 700~K; \texttt{coolant\_pressure=0.2~MPa},
\texttt{initial\_gas\_pressure=0.1~MPa}; \texttt{tend=300~s},
\texttt{dtout=10~s}; tolerances \texttt{(1e-5, 1e-7)}.  300~s is chosen
deliberately: long enough for FRED-OX's $\sim$150--200~s thermal time
constant to reach steady state, short enough that burnup and FGR stay
negligible ($\text{fgrel}, \text{bup} \sim 10^{-4}$--$10^{-22}$, printed
explicitly by the script) --- i.e.\ ``hot steady state, before any serious
irradiation effects'', matching the docstring's own framing.

\paragraph{What to expect.}

\begin{itemize}
  \item \textbf{Cold run:} $T_0=293.15$~K $\to$ peak fuel temperature
        $\approx 1454.83$~K by $t\approx150$~s.
  \item \textbf{Hot run:} flat at $\approx 1454.834$~K for the whole run
        (max deviation from its own $t=0$ value $\approx 1.95\times10^{-4}$~K).
\end{itemize}

\begin{warningbox}[Interpreting ``no serious irradiation effects'']
  The script asserts \texttt{fgrel}/\texttt{bup} stay tiny over the 300~s
  window specifically to prove that the hot-vs-cold temperature comparison
  is a thermal-only effect, not an artifact of irradiation state diverging
  between the two runs.  For a longer run (hours to days), do not assume the
  same holds --- burnup accumulates in real time regardless of hot/cold
  start, so the two runs' irradiation states will diverge once real time
  passes $t=0$, even though they started from the same thermal state.
\end{warningbox}

% ===========================================================================
\section{Example OX-3: HDF5 Output}
\label{sec:ox-hdf5}
% ===========================================================================

\textbf{File:} \texttt{examples/fred\_ox/hdf5\_output/run.py}

\medskip
Extends FRED-ROD's HDF5 example (Section~\ref{sec:hdf5}) with the
FRED-OX-specific fields.  Uses the same \texttt{fred\_input}/\texttt{fred\_output}
pattern:

\begin{lstlisting}
import fred_input as fred
import fred_output as fo

solver.run(tend=86400.0, dtout=1800.0, output_file='results.h5')
solver2.run(tend=7200.0, dtout=3600.0, output_file='debug.h5', all_steps=True)

r = fo.read_results_h5('results.h5')
gpres  = r['gpres']         # (n_steps,) MPa
fggen  = r['fggen']         # (n_steps,) mol
fgrel  = r['fgrel']         # (n_steps,) mol
bup    = r['bup']           # (n_steps,) MWd/kgU
gap_w  = r['gap_width']     # (n_steps,) m
peak_T = r['peak_T_fuel']   # (n_steps,) K
T      = r['T']             # (n_steps, nz, nf+nc) K -- per-layer, unlike ROD's flat array
\end{lstlisting}

\begin{warningbox}[Shape convention]
  Unlike FRED-ROD's single-layer examples, \texttt{r['T']} here is
  genuinely 3-D (\texttt{[n\_steps, nz, nf+nc]}) because \texttt{nz>1}.
  The script uses this to plot an axial temperature profile (fuel
  centreline and clad outer surface vs.\ layer) at the final time step ---
  not meaningful for the single-layer FRED-ROD examples.
\end{warningbox}

Unit conventions for the new fields match Example~OX-1: gas amounts in mol
(convert to cm$^3$ STP the same way), burnup in MWd/kgU, gas pressure in MPa.

% ===========================================================================
\section{Example OX-4: Restart (Checkpoint)}
\label{sec:ox-restart}
% ===========================================================================

\textbf{File:} \texttt{examples/fred\_ox/restart/run.py}

\medskip
Checkpoint-only fault recovery, the FRED-OX analogue of the checkpoint half
of FRED-ROD's Example~7 (Section~\ref{sec:rs}): a single overwritten binary
file for resuming after a crash, with simulation time \emph{preserved} on
reload (as opposed to a snapshot, where time resets to 0 --- see
Example~OX-5).

\textbf{Scenario:} a 3-day (\texttt{tend=3*86400~s}) single-layer MOX/T91
pin (\texttt{nf=4, nc=3, nz=1}), \texttt{pu\_content=0.20},
\texttt{rof0=10200~kg/m$^3$}, cladding \texttt{reference\_density=7750~kg/m$^3$},
power $1.5\times10^8$~W/m$^3$, coolant 620~K, \texttt{coolant\_pressure=0.1~MPa}.

\begin{lstlisting}
solver.set_checkpoint_prefix(ckpt_prefix)
solver.run(tend, dtout)                       # (a) full reference run

s_partial = make_solver()
s_partial.set_checkpoint_prefix(ckpt_prefix)
s_partial.run(86400.0, dtout)                 # (b) crash simulated at day 1

s_restart = make_solver()
s_restart.set_checkpoint_prefix(ckpt_prefix)
s_restart.load_checkpoint(ckpt_file)
print(s_restart.restart_time())               # non-zero, e.g. ~86400 s -- time preserved
s_restart.run(tend, dtout)                    # (c) continue to tend
\end{lstlisting}

The script then verifies that \texttt{peak\_fuel\_temperature()},
\texttt{gas\_pressure()}, \texttt{burnup()}, and \texttt{fg\_released()} from
the restarted run match the full reference run at overlapping output times
(\texttt{rtol=1e-6}, gas pressure \texttt{1e-5}), printing a PASS/FAIL
comparison table --- confirming that irradiation state, not just thermal
state, survives a checkpoint round-trip exactly.

% ===========================================================================
\section{Example OX-5: Restart and Snapshot}
\label{sec:ox-rs}
% ===========================================================================

\textbf{Files:} \texttt{examples/fred\_ox/restart\_snapshot/\{0\_full\_run,1\_checkpoint\_restart,2\_continuation\}.py}

\medskip
The same three-script pattern as FRED-ROD's Example~7 (Section~\ref{sec:rs}),
using the identical 3-day scenario as Example~OX-4, but exercising
\emph{both} persistence mechanisms together and confirming FRED-OX's
irradiation state (burnup, fission-gas inventory) survives each one
correctly:

\begin{description}
  \item[\texttt{0\_full\_run.py}] Reference run.  Enables both
        \texttt{set\_checkpoint\_prefix(prefix)} and
        \texttt{set\_snapshot\_prefix(prefix, snapshot\_times=[t\_half])}.
        Produces \texttt{ox\_checkpoint.chk} (overwritten fault-recovery
        file), \texttt{ox\_frame1.snapshot} (at $t_\text{half}=86400$~s),
        \texttt{ox\_frame2.snapshot} (final state, saved automatically), and
        reference \texttt{.npy} arrays (\texttt{ref\_times/T/gp/bup/fgr.npy}).
  \item[\texttt{1\_checkpoint\_restart.py}] Same pattern as Example~OX-4:
        partial run to day~1 with a checkpoint, reload, \texttt{restart\_time()}
        confirmed non-zero, continue to \texttt{tend}; verified against the
        script-0 reference.
  \item[\texttt{2\_continuation.py}] Snapshot-based two-phase run: Phase~1
        (0~$\to$~$t_\text{half}$) auto-saves a snapshot; Phase~2 calls
        \texttt{s2.load\_snapshot(snap\_file)}, confirms
        \texttt{s2.restart\_time() == 0.0} (time reset) while burnup and
        fission-gas state are preserved from Phase~1, then runs
        $0 \to t_\text{end} - t_\text{half}$.  Trajectories are stitched by
        offsetting Phase~2 times by $t_\text{half}$, and the combined series
        is verified against the script-0 reference.
\end{description}

\begin{infobox}[Key learning point]
  The fields explicitly carried through and checked across every restart
  path here --- \texttt{peak\_fuel\_temperature}, \texttt{gas\_pressure},
  \texttt{burnup}, \texttt{fg\_released} --- are exactly the extra DAE/ODE
  state FRED-OX adds over FRED-ROD.  This example is the most direct proof
  that the gas-plenum-pressure DAE, burnup accumulator, and FGR inventory
  all persist correctly through both time-preserving (checkpoint) and
  time-resetting (snapshot) restarts.
\end{infobox}

% ===========================================================================
\section{Example OX-6: ESFR Pin --- Base Irradiation and Transients}
\label{sec:esfr}
% ===========================================================================

\textbf{Files:} \texttt{examples/fred\_ox/ESFR\_pin/\{base\_irradiation,ULOF\_transient,UTOP\_transient\}/run.py}

\medskip
The most advanced FRED-OX example: a full-length base-irradiation benchmark
followed by two fast, unprotected-transient analyses that pick up from the
irradiated state via a saved snapshot.  ESFR = European Sodium Fast Reactor
(ESFR-SMART); all three scripts share the same pin geometry and materials as
Example~OX-1 (\texttt{nz=20, nf=22, nc=5}, MOX/T91/He).

\subsection{\texttt{base\_irradiation}: 600-day benchmark}

Same setup as Example~OX-1 but run for 600 days
($t_\text{end}=5.184\times10^7$~s, \texttt{dtout=86400~s}), with boundary
conditions given at 5 knot times
($0, 3600, 86400, 2.592\times10^7, 5.184\times10^7$~s) per axial layer.  At
$t=600$~d the legacy reference gives: gas pressure 0.59687~MPa, fission gas
generated 428.871~cm$^3$~STP, released 101.969~cm$^3$~STP (23.776\% FGR),
peak fuel temperature (layer 13) 1826.6~$^\circ$C.  The new call here is:

\begin{lstlisting}
solver.set_snapshot_prefix(prefix)   # auto-saves esfr_pin_base_frame1.snapshot at run end
\end{lstlisting}

This snapshot --- the fully irradiated, 600-day pin state --- is what the
two transient scripts below load.

\subsection{\texttt{ULOF\_transient}: Unprotected Loss of Flow (43~s)}

\begin{lstlisting}
solver = make_solver()
solver.load_snapshot('../base_irradiation/snapshots/esfr_pin_base_frame1.snapshot')
print(solver.restart_time())          # 0.0 -- confirms time reset; irradiation state kept
solver.set_swelling_multiplier(0.0)   # freeze further solid swelling during the fast transient
solver.run(tend=43.0, dtout=0.1)      # 430 output steps
\end{lstlisting}

Models a coolant pump-halving event: 38 boundary-condition time points span
$t=0$ to $43.5$~s (pump halving time 10~s; power drops $\sim$28\% over the
transient as flow-driven reactivity feedback kicks in).  Tolerances
\texttt{(1e-5, 1e-8)}.  \texttt{set\_swelling\_multiplier(0.0)} is set
because 43~s is far too short for any further irradiation-timescale solid
swelling to occur --- only the already-accumulated (loaded-from-snapshot)
swelling and the fast thermal/mechanical response matter here.

\subsection{\texttt{UTOP\_transient}: Unprotected Transient Over-Power (1000~s)}

Same snapshot-loading pattern as ULOF, but models a $+280$~pcm reactivity
insertion: power peaks $\sim$37\% above nominal, then decays, over 52
boundary-condition time points spanning $t=0$ to $1000.4$~s. Also uses
\texttt{set\_swelling\_multiplier(0.0)}; \texttt{dtout=1.0~s} (1000 output
steps); tolerances \texttt{(1e-4, 1e-8)}.

\medskip
Both transient scripts print final gas pressure, peak fuel temperature,
average gap width, and fission gas released, and produce a $2\times2$
subplot grid (\texttt{plots/ulof\_results.png} / \texttt{plots/utop\_results.png})
of these quantities vs.\ time.

\begin{infobox}[Key learning point]
  \texttt{ESFR\_pin} demonstrates the intended FRED-OX workflow for
  transient safety analysis: run (or load from a validated database) a
  long base-irradiation history to build up realistic burnup/gas/swelling
  state, snapshot it, then branch into as many fast transient scenarios as
  needed --- each starting from the same irradiated state at $t=0$ without
  re-running the multi-hundred-day irradiation each time.
\end{infobox}

% ===========================================================================
\section{FRED-OX API Reference}
\label{sec:api-ox}
% ===========================================================================

\subsection{Materials}

\begin{center}
\begin{tabular}{lll}
\toprule
Class & Type & Notes \\
\midrule
\texttt{FredOxMOX(pu\_content, rof0, sto0)} & Fuel     & MOX with base-irradiation correlations \\
\texttt{FredOxT91(reference\_density)}      & Cladding & Ferritic-martensitic, irradiation-aware \\
\texttt{FredOxGapMaterial()}                & Gap      & He/fission-gas mixture, burnup-coupled \\
\bottomrule
\end{tabular}
\end{center}

\subsection{Solver methods: \texttt{FredOxSolver} --- additions over \texttt{FredRodSolver}}

\begin{center}
\begin{tabular}{ll}
\toprule
Method & Description \\
\midrule
\texttt{set\_power\_density\_history\_per\_layer(t, qqv)} & Per-layer volumetric heat [W/m$^3$] \\
\texttt{set\_coolant\_temperature\_per\_layer(t, T)}       & Per-layer coolant BC [K] \\
\texttt{set\_plenum\_temperature\_history(t, T)}           & Gas-plenum temperature [K] \\
\texttt{set\_swelling\_multiplier(f)}                      & Scale solid fission-product swelling \\
\texttt{set\_hot\_start(bool)}                              & Pre-march with irradiation frozen (Section~\ref{sec:hot-ox}) \\
\midrule
\texttt{gas\_pressure()}              & Gas-plenum pressure [MPa] (extra DAE state) \\
\texttt{fg\_generated()}              & Total fission gas generated [mol] \\
\texttt{fg\_released()}               & Total fission gas released [mol] \\
\texttt{burnup()}                     & Rod burnup [MWd/kgU] \\
\bottomrule
\end{tabular}
\end{center}

All other geometry attributes, materials-interface concepts, and solver
methods (temperatures, gap width, checkpoint/snapshot, HDF5 output) are
identical to FRED-ROD --- see Section~\ref{sec:api-rod}.

% ===========================================================================
\section{FRED-OX Troubleshooting}
\label{sec:trouble-ox}
% ===========================================================================

\begin{description}[leftmargin=0pt]
  \item[Fission gas release seems to jump discontinuously across a restart.]
    Confirm you used \texttt{load\_checkpoint} (time-preserving) and not
    \texttt{load\_snapshot} (time-resetting) if you intended to continue the
    \emph{same} irradiation history --- see Section~\ref{sec:ox-restart}.
    Loading a snapshot deliberately resets simulation time to 0 while
    keeping burnup/gas state, which is correct for starting a new transient
    from an irradiated pin, but wrong if you wanted a literal continuation.

  \item[\texttt{set\_hot\_start(True)} gives an unexpectedly large burnup at $t=0$.]
    This should not happen --- irradiation physics is frozen during the
    pre-march (Section~\ref{sec:hot-ox}).  If you see nonzero burnup growth
    at $t=0$, check you are calling \texttt{set\_hot\_start} before
    \texttt{run()}, not after building a solver you have already partially
    run.

  \item[Per-layer boundary condition array shape errors.]
    \texttt{set\_power\_density\_history\_per\_layer} and
    \texttt{set\_coolant\_temperature\_per\_layer} expect a 2-D array shaped
    \texttt{(n\_time\_points, nz)} --- one column per axial layer, in the
    same top-to-bottom order as \texttt{g.dz0}.  A single-layer (\texttt{nz=1})
    example still needs shape \texttt{(n, 1)}, not a flat 1-D array.

  \item[Legacy comparison mismatch after changing \texttt{set\_swelling\_multiplier}.]
    The legacy reference values quoted in Examples~OX-1 and OX-6 assume the
    default multiplier used in those scripts (0.8 for \texttt{base\_irradiation},
    unset/1.0 elsewhere).  Changing it is a deliberate tuning choice and will
    shift gap width, gas pressure, and hence FGR relative to the quoted
    reference numbers.

  \item[See also] FRED-ROD's troubleshooting list (Section~\ref{sec:trouble-rod})
    for issues common to both apps (\texttt{IDACalcIC} failures, gap-closure
    setup, HDF5/\texttt{h5py} availability, Python module path issues).
\end{description}

% ===========================================================================
\section{Beyond FRED-ROD and FRED-OX: FRED-M-Na}
\label{sec:mna-pointer}
% ===========================================================================

FRED-M-Na (metallic U-Pu-Zr fuel in sodium) is built on the same platform
layer as FRED-ROD and FRED-OX but is not covered in depth by this manual.
Its examples live under \texttt{examples/fred\_m\_na/} and follow the same
general pattern (geometry $\to$ materials $\to$ solver $\to$ run), with its
own physics documented in \texttt{doc/physics\_fred\_m\_na.md} and
\texttt{doc/theory\_manual/fred\_m\_na.tex}.

\begin{center}
\begin{tabular}{ll}
\toprule
Example & What it shows \\
\midrule
\texttt{simple\_irradiation}         & First FRED-M-Na run: single-pin base irradiation \\
\texttt{hot\_start\_demo}            & Cold vs hot start (same concept as Sections~\ref{sec:hot-rod}, \ref{sec:hot-ox}) \\
\texttt{restart\_snapshot}           & Checkpoint and snapshot persistence (same pattern as Section~\ref{sec:ox-rs}) \\
\texttt{timpano\_SAS\_benchmark}     & Multi-year benchmark against legacy FRED-M / SAS4A \\
\bottomrule
\end{tabular}
\end{center}

\begin{infobox}[Where to start with FRED-M-Na]
  If you are already comfortable with FRED-ROD and FRED-OX from this
  manual, \texttt{simple\_irradiation} $\to$ \texttt{hot\_start\_demo} $\to$
  \texttt{restart\_snapshot} follows the same learning progression as
  Parts~\ref{part:rod}--\ref{part:ox}, before moving on to the full
  \texttt{timpano\_SAS\_benchmark} legacy-validation case.
\end{infobox}

% ===========================================================================
\end{document}
% ===========================================================================
